\section*{LỜI NÓI ĐẦU}
\thispagestyle{empty}
Trong những năm gần đây, công nghệ phục hồi chức năng đã có những bước phát triển vượt bậc, đóng góp quan trọng trong việc cải thiện chất lượng cuộc sống của những người bị tổn thương vận động do tai nạn, đột quỵ hoặc các bệnh lý khác. Việc ứng dụng các thiết bị cơ khí kết hợp điều khiển điện tử thông minh không chỉ giúp tăng hiệu quả luyện tập mà còn giảm thiểu sự phụ thuộc vào sức lao động của nhân viên y tế.

Đồ án “Thiết kế và chế tạo thiết bị phục hồi chức năng cánh tay” được thực hiện nhằm mục tiêu xây dựng một hệ thống đơn giản, hiệu quả, có khả năng mô phỏng các chuyển động nâng hạ và gập duỗi của cánh tay, hỗ trợ quá trình phục hồi vận động cho bệnh nhân. Thiết bị sử dụng động cơ bước làm cơ cấu truyền động kết hợp bộ điều khiển ESP8266 với các cảm biến đo vị trí và an toàn nhằm đảm bảo sự chính xác và an toàn trong quá trình vận hành.

Đồ án không chỉ mang ý nghĩa thực tiễn trong lĩnh vực y tế và phục hồi chức năng mà còn giúp tác giả nâng cao kỹ năng thiết kế cơ khí, lập trình điều khiển nhúng và tích hợp hệ thống đa lĩnh vực. Qua đó, góp phần phát triển năng lực chuyên môn, đồng thời đáp ứng nhu cầu thực tiễn trong ứng dụng công nghệ hỗ trợ phục hồi sức khỏe.

Tôi xin chân thành cảm ơn sự hướng dẫn tận tình của thầy NGUYỄN PHAN KIÊN, sự hỗ trợ của bạn bè và gia đình trong suốt quá trình thực hiện đề tài này.
\cleardoublepage